%% Owner: Víctor Martín Castilla
\section{Introduction}
\label{sec:introduction}

The short project asks for the digital transformation of a company's
infrastructure into containerised microservices
\parencite{llusa2026project}. The company runs four groups of services,
each deployed as its own Docker stack: shared basic services, a core
application, an observability platform and a Moodle learning management
system. Everything must run on a Docker Swarm of at least three nodes,
be described by Compose files, and scale to several containers, except
for cron jobs and database engines that cannot.

Deploying may only take two kinds of step: building images and deploying
the Compose files. Two Compose files are required per stack: one for a
local environment without TLS or mapped volumes, and one with the full
production configuration.

The team has three members rather than the five the statement assumes,
so the work was split by stack, with one owner per stack and every
change reviewed through a pull request (\cref{sec:team}).

\Cref{sec:architecture} describes the services and how they map to
stacks. \Cref{sec:infrastructure} explains where the cluster runs and how
its virtual machines were created and secured, and \cref{sec:swarm} how
the Swarm was built on them. Sections~\ref{sec:stack-basic}
to~\ref{sec:stack-stats} cover each stack, \cref{sec:data} the data and
its backups, \cref{sec:hadr} high availability and disaster recovery, and
\cref{sec:deployment} the deployment procedure.
